\begin{mistake}{2026-08-02}{03}

\begin{problemcontent}
设曲线 \(y = y(x)\) 有二阶连续导数，\(y''(x) > 0\)，其上任意一点 \((x, y)\) 处的曲率 \(K = \frac{1}{2y^2 \cos \alpha}\)（\(\cos \alpha > 0\)），其中 \(\alpha\) 为该曲线在相应点处的切线的倾角，且该曲线在点 \((1, 1)\) 处取得极小值，求曲线 \(y = y(x)\)。
\end{problemcontent}

\begin{solutioncontent}
由导数几何意义，\(\tan \alpha = y'\)，又因为 \(\cos \alpha > 0\)，所以 \(\cos \alpha = \frac{1}{\sqrt{1 + (y')^2}}\)。
因为 \(y''(x) > 0\)，曲率 \(K = \frac{y''}{(1 + (y')^2)^{\frac{3}{2}}}\)。代入已知条件得：
\[
\frac{y''}{(1 + (y')^2)^{\frac{3}{2}}} = \frac{\sqrt{1 + (y')^2}}{2y^2} \implies 2y^2 y'' = (1 + (y')^2)^2
\]
此为缺 \(x\) 型微分方程。令 \(y' = p\)，则 \(y'' = p \frac{dp}{dy}\)，代入方程得：
\[
2y^2 p \frac{dp}{dy} = (1 + p^2)^2 \implies \frac{2p}{(1 + p^2)^2} dp = \frac{1}{y^2} dy
\]
两边积分得：
\[
-\frac{1}{1 + p^2} = -\frac{1}{y} + C_1
\]
曲线在 \((1,1)\) 处取得极小值，故 \(y(1) = 1, y'(1) = p(1) = 0\)，代入得 \(C_1 = 0\)。
化简得 \(1 + (y')^2 = y\)，即 \((y')^2 = y - 1\)。

对等式两边关于 \(x\) 求导：
\[
2y' y'' = y'
\]
因 \(y'' > 0\)，\(y'\) 不恒为 \(0\)，约去 \(y'\) 得 \(2y'' = 1\)，即 \(y'' = \frac{1}{2}\)。
积分得 \(y' = \frac{1}{2}x + C_2\)。代入 \(y'(1) = 0\) 得 \(C_2 = -\frac{1}{2}\)，即 \(y' = \frac{1}{2}(x - 1)\)。
再积分得 \(y = \frac{1}{4}(x - 1)^2 + C_3\)。代入 \(y(1) = 1\) 得 \(C_3 = 1\)。
故所求曲线方程为：
\[
y = 1 + \frac{1}{4}(x - 1)^2
\]
\end{solutioncontent}

\mistakeknowledge{
\begin{itemize}
  \item \textbf{核心公式}：曲率 \(K = \frac{|y''|}{(1 + (y')^2)^{\frac{3}{2}}}\)，切线倾角与导数关系 \(\sec^2 \alpha = 1 + (y')^2\)。
  \item \textbf{易错点}：曲率公式中的绝对值易漏，需利用题目条件 \(y'' > 0\) 去掉绝对值。
  \item \textbf{计算技巧}：得到 \((y')^2 = y - 1\) 后，直接开方需分情况讨论，两边直接对 \(x\) 再次求导可巧妙避开符号讨论。
\end{itemize}}

\end{mistake}
